%% ============================================================
\section{The \coeaudit{}}
\label{sec:coeaudit}
%% ============================================================


\begin{figure}[t]
\centering
\includegraphics[width=0.9\textwidth]{figures/06-forensic-audit-overview.jpg}
\caption{\coeaudit{} overview. An adapter normalizes each system's deliverables (\texttt{paper.tex}, solution code, \texttt{references.bib}) into a common artifact bundle, on which four integrity checks run independently: \textbf{I1 Score Verification} re-runs the solution on the golden evaluator and compares to the extracted paper score within an adaptive tolerance (match/mismatch); \textbf{I2 Specification Violation} uses majority-vote LLM judgment over solution code, evaluator, and task spec (clean/flagged); \textbf{I3 Reference Verification} resolves each bib entry via academic APIs with LLM disambiguation of near-misses (verified/hallucinated); and \textbf{I4 Method--Code Alignment} uses majority-vote LLM judgment of method described in the paper vs.\ solution code (aligned/misaligned).}
\label{fig:forensic_audit}
\end{figure}



\emph{\coeaudit{}} is a post-hoc audit that checks whether claims in a completed paper are supported by the underlying artifacts---code, evaluator outputs, and bibliography (Figure~\ref{fig:forensic_audit}).
It comprises four integrity checks, each targeting a tractably verifiable claim type.
While the CoE framework can support other evaluation forms---real-time verification during paper production or broader claim coverage---those are outside the scope of this work.
An adapter unifies each system's deliverables (paper, solution code, references) into a common artifact bundle, then the four checks run independently, each targeting a specific way a claim can lose its grounding:
\label{sec:coeaudit:modes}


\paragraph{Score Verification (I1).}
\label{sec:coeaudit:i1}
The paper's reported score is extracted by LLMs from both \TeX{} and PDF files, then compared against scores obtained by re-running the submitted solution on the golden evaluator.
A paper passes if the two match within an adaptive tolerance that accounts for evaluator noise.

\paragraph{Specification Violation (I2).}
\label{sec:coeaudit:i2}
Specification violations occur when solution code breaks task rules---for example, reverse-engineering the evaluator's scoring logic, or hardcoding answers for known test cases.
The generating agent optimizes for the score rather than genuinely solving the problem the task poses.
LLMs inspect the solution code against the golden evaluator and task specification to detect such violations, with majority vote across multiple runs.

\paragraph{Reference Verification (I3).}
\label{sec:coeaudit:i3}
Each bibliography entry is resolved by querying multiple academic APIs (Semantic Scholar, arXiv, OpenAlex, CrossRef) using arXiv ID, DOI, and title.
An LLM cross-checks the full bib entry against returned records to catch near-misses and citation gaming (e.g., a real DOI attached to a fabricated description).
Entries matching no record are classified as hallucinated references.

\paragraph{Method-Code Alignment (I4).}
\label{sec:coeaudit:i4}
An LLM reads the paper's method section and the solution code side by side, then judges whether the paper faithfully describes what the code does.
Acceptable simplification (e.g., omitting implementation details) is treated as aligned; only cases where the paper describes a fundamentally different algorithm count as misaligned.
We conduct multiple independent runs with majority vote to reduce LLM judgment noise.

\paragraph{Native claim provenance.}
The four checks above are \emph{forensic}: they operate on submitted artifacts alone and apply identically to every system.
For systems that emit structured provenance at write-time---linking each claim to a specific source record---an additional \emph{native} check becomes possible: the \textbf{numerical Claim Provenance Rate (CPR)}, which measures the fraction of quantitative claims in the paper that trace to a matching entry in the experimental log.
We report this check for \sys{}, the only system in our evaluation that produces such provenance records (\S\ref{sec:exp:native_cpr}).
